\SourceLeaf{58}
\fancyhead[L]{\small ANZEIGEN EIGENER SCHRIFTEN}
\pdfbookmark[1]{Theoria motus corporum coelestium. Anzeige von 1809}{anzeige-theoria1809}
\begin{center}\rule{.75\textwidth}{.3pt}\\
Göttingische gelehrte Anzeigen. 17. Juni 1809.\\
\rule{.75\textwidth}{.3pt}\end{center}
\textit{Theoria motus corporum coelestium in sectionibus conicis Solem ambientium.} Von \textsc{Carl Frid. Gauß.} Hamburg, 1809. Im Verlag von \textit{Frid. Perthes} und \textit{J. H. Besser.} XII Seiten Vorrede, 228 Seiten Text und 20 Seiten Tabellen, dazu eine Kupfertafel. Großquart.

Mit dem Teil der theoretischen Astronomie, der die parabolische und elliptische Bewegung der Himmelskörper zum Gegenstand hat, haben sich bekanntlich viele Autoren beschäftigt, darunter sogar Mathematiker ersten Ranges in eigenen Werken. Nach solchen Vorgängern konnte nur ein gewichtiger Anlass eine neue Bearbeitung dieses Gebiets rechtfertigen. Den eigentlichen Astronomen ist dieser Anlass hinreichend bekannt. Nur jenen Freunden der Himmelskunde, die daraus keine Hauptbeschäftigung machen und vielleicht in den letzten verhängnisvollen Jahren über den Angelegenheiten der Erde die Angelegenheiten des Himmels aus den Augen verloren haben, wollen wir kurz die Umstände in Erinnerung rufen, die unmittelbar zu diesem Werk geführt haben. Die Aufgabe besteht darin, aus den nur mäßig lange von der Erde aus beobachteten Bewegungen eines Himmelskörpers seine Bahn mit hinreichender Genauigkeit zu bestimmen, wenn man von ihm nur weiß, dass er sich nach den Gesetzen \textsc{Kepler}s in einem Kegelschnitt um die Sonne bewegt. Bisher war diese Aufgabe eigentlich noch nie
\SourceLeaf{59}ernsthaft bearbeitet worden. Zwar gibt es einige Untersuchungen zu diesem Problem. Ohne den Scharfsinn und die analytische Kunst zu verkennen, von denen einige dieser Untersuchungen Spuren zeigen, muss man den vorgeschlagenen Methoden doch alles absprechen, was zur praktischen Brauchbarkeit nötig ist. Sie liefern statt der größtmöglichen Genauigkeit bestenfalls eine grobe Annäherung oder bieten statt eines flexiblen, auf die praktische Anwendung zugeschnittenen Rechenverfahrens nur einen verworrenen Haufen unausgearbeiteter Formeln, die selbst den unermüdlichsten Rechner abschrecken. Oder sie werden schließlich schon durch weit kleinere Fehler als die bei Beobachtungen unvermeidlichen völlig unbrauchbar, obwohl sie auf Beobachtungen anwendbar sein müssten, wie sie der heutige Stand der praktischen Astronomie erlaubt. Dass ein schon an sich zweifellos so interessantes Problem vernachlässigt wurde, scheint zum Teil daran zu liegen, dass die Mathematiker, die sich damit beschäftigten, mit den Möglichkeiten und Bedürfnissen der Praxis nicht genügend vertraut waren. Vor allem lag es aber wohl daran, dass die Geschichte der Astronomie noch keinen Fall kannte, in dem eine angemessene Lösung der Aufgabe dringend nötig und ihr Nutzen deutlich spürbar gewesen wäre. Als nämlich \textsc{Kepler} nach der Entdeckung seiner Gesetze daran ging, die Abmessungen der Bahnen der damals bekannten Planeten zu bestimmen, standen ihm neben den bereits sehr genau bekannten mittleren Bewegungen ein reicher Schatz guter Beobachtungen \textsc{Tycho}s aus vielen Jahren zur Verfügung. Daraus brauchte er jeweils nur auszuwählen, was er für seine schönen, aber auf Sonderfälle zugeschnittenen und vergleichsweise einfachen Methoden benötigte. Dieselben Hilfsmittel, ja noch bessere, hatten \textsc{Kepler}s Nachfolger. Mit den Fortschritten der Beobachtungskunst machten sie sich daran, die schon recht gut bekannten Elemente der Planetenbahnen noch genauer zu bestimmen.

Anders lagen die Dinge bei den Kometen, in deren Bewegung man nach \textsc{Newton}s Entdeckungen nur einen Sonderfall der allgemeinen Gesetze \textsc{Kepler}s erkannte. Diese Himmelskörper sind gewöhnlich nur kurze Zeit sichtbar. Wie ihre Bewegungen erscheinen, hängt davon ab, welche Stellungen die Erde während dieser Zeit gerade einnimmt. Der Mathematiker, der aus den Erscheinungen die Bahnen bestimmen will,
\SourceLeaf{60}hat unter den Beobachtungen wenig Auswahl. Er muss sie nehmen, wie er sie vorfindet, wie der Zufall sie ihm gegeben hat. Nur äußerst selten stehen ihm Beobachtungen zur Verfügung, die unter den Umständen angestellt wurden, die er sich für die Anwendung spezieller Methoden wünschen würde. Die Aufgabe, aus einigen Beobachtungen eines Kometen seine parabolische Bahn zu bestimmen, war daher ebenso wichtig wie schwierig. \textsc{Newton} selbst erkannte die Schwierigkeiten an und wusste sie zu überwinden. Mit welchem Eifer und mit welchem Erfolg man sich seit \textsc{Newton} bis in unsere Zeit mit dieser Aufgabe beschäftigt hat, ist hinreichend bekannt. Zwischen dieser Aufgabe und der oben von uns erwähnten besteht aber ein wesentlicher Unterschied. Bei jener wird die Art des Kegelschnitts \textit{vorgeschrieben}. Bei dieser soll sie aus der unendlichen Vielfalt möglicher Kegelschnitte, von denen die Parabel nur eine Art ist, allein mithilfe der Beobachtungen selbst und ohne hypothetische Voraussetzungen ermittelt werden. Verständlicherweise vereinfacht diese Einschränkung jene Aufgabe erheblich. Bei den meisten Kometen durfte und musste man sich diese Vereinfachung gestatten, weil die gewöhnlich recht geringe Genauigkeit der Beobachtungen und ihre kurze Dauer kaum je ausreichen, um eine Abweichung von der Parabel nachzuweisen und ihre Größe zu bestimmen. Bei einigen Kometen hat man diese Bestimmung zwar tatsächlich versucht, aber, und das ist wesentlich, stets erst, nachdem man bereits eine parabolische Bahn berechnet hatte, die als Annäherung an die zu bestimmende Ellipse oder Hyperbel diente. So blieb unser allgemeineres Problem auch hier gewissermaßen entbehrlich.

Auch die Entdeckung eines neuen dauerhaft vorhandenen Himmelskörpers im Jahr 1781 ließ den Bedarf an einer Lösung dieses Problems noch nicht spürbar werden. Nachdem sich eine parabolische Bahn als unzureichend erwiesen hatte, versuchte man es mit einem Kreis. Weil die Exzentrizität der Planetenbahn zufällig recht klein war, stellte er die Bewegung des Planeten einige Jahre lang leidlich genau dar. Angesichts seiner sehr langsamen Bewegung, seines geringen Abstands von der Ekliptik und seiner noch recht großen Helligkeit bestand keine \textit{Gefahr im Verzug}. Man fand ihn von einem Jahr zum anderen mühelos wieder und konnte mit der Bestimmung der Abweichung vom Kreis oder der wahren Ellipse warten, bis man die Beobachtungen nach Belieben auswählen konnte.

\SourceLeaf{61}Ganz anders verhielt es sich mit der im Jahr 1801 entdeckten Ceres. Dieser Himmelskörper erscheint nur als Sternchen achter Größe. Aus der zahllosen Schar ganz ähnlicher Fixsterne lässt er sich nur mühsam und bei genauer Kenntnis seiner jeweiligen Position herausfinden. Der Entdecker hatte ihn nur 41 Tage lang beobachtet. Als die Entdeckung im übrigen Europa bekannt wurde, war er bereits in den Sonnenstrahlen verschwunden und sollte erst nach einem Jahr in einer ganz anderen Himmelsgegend wieder sichtbar werden. Nun kam es darauf an, den Ort, an dem man ihn wieder suchen musste, möglichst genau vorherzusagen. Die Vorhersage musste allein auf den wenigen vorhandenen Beobachtungen und einer strengen Rechnung beruhen, ohne unsichere Hypothesen. Mehrere Astronomen versuchten es mit der einfachsten Hypothese, einer Kreisbahn. Die Beobachtungen ließen sich damit freilich nur unvollkommen in Einklang bringen. Ob die Hypothese zur Wiederauffindung ausreichen würde, blieb also zumindest höchst fraglich. Tatsächlich bestätigte sich später, dass die Kreishypothese schon am Ende des Jahres 1801 um elf Grad vom wahren Ort des Planeten abwich. Die Astronomen, die das Glück hatten, ihn wiederzufinden, erklärten selbst, dass die Wiederauffindung nach einer so fehlerhaften Hypothese unmöglich gewesen wäre.

Dem Verfasser des vorliegenden Werks waren im Sommer 1801 bei einer ganz anderen Beschäftigung einige Ideen gekommen, die ihn zu einer Lösung des erwähnten allgemeinen Problems führen zu können schienen. Zu einer anderen Zeit hätte er diese Ideen, die für ihn zunächst nur einen theoretischen Reiz hatten, vielleicht nicht gleich weiterverfolgt und ausgearbeitet. Doch gerade zu diesem Zeitpunkt hatte \textsc{Piazzis} Entdeckung die allgemeine Aufmerksamkeit gefesselt und waren seine Beobachtungen soeben veröffentlicht worden. Der Verfasser konnte es sich daher nicht versagen, daran die praktische Anwendbarkeit seiner Ideen zu prüfen. Der Erfolg dieser Arbeit ist bekannt. Die bis dahin ungeahnte Möglichkeit war damit aufs Schönste erwiesen, aus einer kurzen Reihe von Beobachtungen eines Planeten seine Bahn schon sehr nahe an der wahren Bahn und weit genauer zu bestimmen, als für seine Wiederauffindung nach längerer Zeit nötig war. Die angewandte Methode hatte sich bewährt. Hätten an ihrer allgemeinen Brauchbarkeit noch Zweifel bestehen können, so wurden diese durch ebenso
